% MAT670 -- Interlude B (end of Part II): lattices, quadratic forms and the geometry of numbers.
% Entirely new material (not in the handwritten notes). Labels prefixed NT:.

\chapter[Interlude B: Number theory]{Interlude B: Lattices, quadratic forms and the geometry of numbers}
\label{NT:ch}

\begin{addedblock}
This interlude is new: it is not in the handwritten notes. It starts from the title of Blichfeldt's
1929 paper quoted in Lecture~2, \emph{``The minimum value of quadratic forms, and the closest packing of
spheres''}, and follows the first half of that title into number theory. To keep it readable it is set
in black; all of it should be read as added material.
\end{addedblock}

Blichfeldt's title puts two problems side by side, and in Lectures~1--6 we only looked at the second
one. This interlude is about the first. We shall see that
\begin{itemize}
\item a lattice \emph{is} a positive definite quadratic form, so the lattice packing problem is the
  question of how large the minimum of a form of given determinant can be (Hermite's problem);
\item counting lattice points in balls is the classical problem of counting representations of integers
  by forms, and the leading term of the count is the covolume, which is what density measures;
\item Blichfeldt's averaging trick from Lecture~2 has a famous cousin, Blichfeldt's lemma of 1914, which
  gives Minkowski's convex body theorem, and through it Fermat's and Lagrange's theorems on sums of squares;
\item the generating function of the lattice point counts, the \emph{theta series}, is a modular form.
  Poisson summation turns this into strong restrictions on very good lattices, and the same Poisson
  summation is the engine of the Cohn--Elkies bound behind Viazovska's theorem.
\end{itemize}
Throughout, $\Lat\subset\R^n$ is a full-rank lattice, $\covol(\Lat)=\abs{\det\Lat}$ as in Lecture~4,
$\lambda_1(\Lat)$ is the length of a shortest nonzero vector, $V_n=\vol(\Ball^n)$, and
\[
  \dens(\Lat)=\frac{V_n\,(\lambda_1/2)^n}{\covol(\Lat)}
\]
is the density of the associated packing (Lecture~4). To avoid a clash with the centre density
$\cent$, the theta series will be written $\Theta_\Lat$.

%======================================================================
\section{Lattices are quadratic forms}
\label{NT:sec:forms}

Let $B$ be a basis matrix of $\Lat$, with the basis vectors $b_1,\dots,b_n$ as its columns. A lattice
vector is $Bx$ with $x\in\Z^n$, and its squared length is
\[
  \abs{Bx}^2=x^TB^TBx=x^TGx=:Q(x),\qquad G=B^TB=\big(\ip{b_i}{b_j}\big)_{i,j}.
\]
The \emph{Gram matrix} $G$ is symmetric and positive definite, and $\det G=\covol(\Lat)^2$. Conversely,
every positive definite symmetric $G$ factors as $G=B^TB$ (Cholesky), and $B$ is unique up to
$B\mapsto UB$ with $U\in\Orth(n)$. Changing the basis replaces $B$ by $BK$ with $K\in\GL_n(\Z)$
(Lecture~4), hence $G$ by $K^TGK$. So:
\[
  \frac{\{\text{lattices in }\R^n\}}{\text{rotations}}
  \;\;\longleftrightarrow\;\;
  \frac{\{\text{positive definite }n\times n\text{ symmetric matrices}\}}{G\sim K^TGK,\ K\in\GL_n(\Z)} .
\]
Under this dictionary
\[
  \lambda_1(\Lat)^2=\min_{x\in\Z^n\setminus0}Q(x)=:\min Q,\qquad \covol(\Lat)=\sqrt{\det G},
\]
and the number of lattice vectors of squared length $m$ is the number of \emph{representations} of $m$
by $Q$. Integral lattices (Lecture~4) are the forms with $G\in M_n(\Z)$. For example $\Z^n$ is the sum
of squares $x_1^2+\dots+x_n^2$, and the hexagonal lattice with basis $(1,0)$,
$(\tfrac12,\tfrac{\sqrt3}2)$ is
\[
  Q_{\mathrm{hex}}(x,y)=x^2+xy+y^2,\qquad G=\begin{pmatrix}1&\tfrac12\\ \tfrac12&1\end{pmatrix},
\]
while the root lattice $A_2$ (Lecture~6, minimal norm $2$) is $2Q_{\mathrm{hex}}$.

\begin{definition}[Hermite's constant]\label{NT:def:hermite}
For a positive definite form put $\gamma(Q)=\min Q/(\det G)^{1/n}$; this is invariant under scaling and
$\GL_n(\Z)$. \emph{Hermite's constant} is $\gamma_n=\sup_Q\gamma(Q)$, the supremum over all positive
definite $n$-ary forms. Equivalently, $\gamma_n=\sup_\Lat\lambda_1(\Lat)^2/\covol(\Lat)^{2/n}$.
\end{definition}

Comparing with the formula for $\dens(\Lat)$ gives
\begin{equation}\label{NT:eq:hermite-density}
  \dens(\Lat)=\frac{V_n}{2^n}\,\gamma(\Lat)^{n/2},
  \qquad\text{so}\qquad
  \sup_\Lat\dens(\Lat)=\frac{V_n}{2^n}\,\gamma_n^{n/2}.
\end{equation}
\emph{Finding the densest lattice packing is the same problem as computing Hermite's constant.} The
known values are below. They are attained exactly by the lattices of Lectures~4--6, and the reader
can recompute each from minimal norm $2$ and the determinants $\det G=n+1$ for $A_n$, $4$ for $D_n$, and
$3,2,1$ for $E_6,E_7,E_8$. For the Leech lattice the minimal norm is $4$ and $\det G=1$.
\begin{center}
\renewcommand{\arraystretch}{1.2}
\begin{tabular}{lccccccccc}
\toprule
$n$ & 1 & 2 & 3 & 4 & 5 & 6 & 7 & 8 & 24\\
\midrule
$\gamma_n^n$ & $1$ & $\tfrac43$ & $2$ & $4$ & $8$ & $\tfrac{64}3$ & $64$ & $2^8$ & $4^{24}$\\
optimal lattice & $\Z$ & $A_2$ & $A_3$ & $D_4$ & $D_5$ & $E_6$ & $E_7$ & $E_8$ & $\Lambda_{24}$\\
density & $1$ & $.9069$ & $.7405$ & $.6169$ & $.4653$ & $.3729$ & $.2953$ & $.2537$ & $.00193$\\
\bottomrule
\end{tabular}
\end{center}
(Lagrange for $n=2$; Gauss for $n=3$; Korkine--Zolotarev for $n=4,5$; Blichfeldt for $n=6,7,8$;
Cohn--Kumar for $n=24$. See the Gallery for references.)

\subsection*{Binary forms: Lagrange's reduction}

For $n=2$ write $Q(x,y)=ax^2+bxy+cy^2$, so $G=\left(\begin{smallmatrix}a&b/2\\b/2&c\end{smallmatrix}\right)$
and $\det G=ac-b^2/4$. Following Lagrange (1773) and Gauss (\emph{Disquisitiones}, 1801), call $Q$
\emph{reduced} if
\[
  \abs b\le a\le c .
\]

\begin{proposition}[Lagrange reduction]\label{NT:prop:reduction}
Every positive definite binary form is equivalent under $\SL_2(\Z)$ to a reduced form, and if $Q$ is
reduced then $\min Q=a$.
\end{proposition}

\begin{proof}
Let $v_1$ be a shortest nonzero vector of the lattice $\Lat$ of $Q$. It is primitive, since
$v_1=kw$ with $k\ge2$ would make $w$ shorter. So $v_1$ extends to a basis $(v_1,w)$, and every other
completion of $v_1$ to a basis is $\pm w+kv_1$, $k\in\Z$. Among the vectors $w+kv_1$ choose a shortest
one, $v_2$, and replace $v_2$ by $-v_2$ if needed to make the determinant positive. In the basis
$(v_1,v_2)$ we have $a=\abs{v_1}^2=\min Q$, $c=\abs{v_2}^2\ge a$, and $b=2\ip{v_1}{v_2}$. Minimality
of $v_2$ gives $c\le\abs{v_2\pm v_1}^2=c\pm b+a$, that is, $\abs b\le a$. For the second statement:
if $\abs b\le a\le c$ and $\abs x\ge\abs y>0$, then
$Q(x,y)\ge a x^2-a\abs{x}\abs{y}+cy^2\ge cy^2\ge a$; the case $\abs y\ge\abs x>0$ is symmetric
(use $\abs b\le c$), and $Q(x,0)=ax^2\ge a$, $Q(0,y)=cy^2\ge a$.
\end{proof}

\begin{theorem}[Lagrange; $\gamma_2=2/\sqrt3$]\label{NT:thm:gamma2}
For every positive definite binary form, $\min Q\le\frac{2}{\sqrt3}\sqrt{\det G}$, with equality only
for forms equivalent to a multiple of $x^2+xy+y^2$. Consequently every lattice packing of disks has
density at most $\pi/\sqrt{12}$, with equality only for the hexagonal lattice.
\end{theorem}

\begin{proof}
Take $Q$ reduced. Then $\det G=ac-\tfrac{b^2}4\ge a^2-\tfrac{a^2}4=\tfrac34a^2$, so
$\min Q=a\le\sqrt{4/3}\sqrt{\det G}$. Equality forces $c=a$ and $\abs b=a$, i.e.\
$Q=a(x^2\pm xy+y^2)$, and $x^2-xy+y^2$ is equivalent to $x^2+xy+y^2$ via $y\mapsto-y$. The density
statement is \eqref{NT:eq:hermite-density}: $\frac{\pi}{4}\cdot\frac{2}{\sqrt3}=\frac{\pi}{2\sqrt3}$.
\end{proof}

This is Theorem~\ref{L1:thm:lattice2} of Lecture~1 again. In fact the proof there, which stacked the
rows $\Z v$ at height $h$ and slid the next row so that one of its points sits over the gap, is
exactly the choice of $v_1$ and $v_2$ above. That proof generalizes to all $n$ as Hermite's inequality
(Exercise~\ref{NT:ex:hermite}).

\subsection*{The picture: the modular figure}

Rotate and scale so that $v_1=1\in\C$, and write $v_2=\tau$ with $\operatorname{Im}\tau>0$. Then
\[
  Q(x,y)=\abs{v_1}^2\abs{x+y\tau}^2=a\big(x^2+2\operatorname{Re}\tau\,xy+\abs\tau^2y^2\big),
\]
so $b/a=2\operatorname{Re}\tau$ and $c/a=\abs\tau^2$. The reduction conditions become
\[
  \abs{\operatorname{Re}\tau}\le\tfrac12,\qquad \abs\tau\ge1 ,
\]
which describe the classical fundamental domain $\mathcal F$ of $\SL_2(\Z)$ acting on the upper half
plane $\mathbb H$ by M\"obius transformations. A change of basis in $\SL_2(\Z)$ acts on $\tau=v_2/v_1$
exactly by $\tau\mapsto\frac{a\tau+b}{c\tau+d}$. So the space of planar lattices up to rotation and
scaling is $\SL_2(\Z)\backslash\mathbb H$, and $\mathcal F$ is a map of it (Figure~\ref{NT:fig:modular}).
Allowing $\GL_2(\Z)$ also identifies $\tau$ with $-\bar\tau$, which folds $\mathcal F$ in half.

In these coordinates the density is very simple. The covolume is $\abs{v_1}^2\operatorname{Im}\tau$ and,
for $\tau\in\mathcal F$, the minimum is $\abs{v_1}^2$, so
\[
  \dens(\tau)=\frac{\pi\abs{v_1}^2/4}{\abs{v_1}^2\operatorname{Im}\tau}=\frac{\pi}{4\operatorname{Im}\tau}
  \qquad(\tau\in\mathcal F).
\]
The level sets of the density are horizontal lines. The densest lattice is the \emph{lowest point} of
$\mathcal F$, the corner $\rho=e^{\pi i/3}$ with $\operatorname{Im}\rho=\frac{\sqrt3}{2}$: the hexagonal
lattice. The square lattice is $\tau=i$, with $\dens=\pi/4$. Lattices high up in the cusp are long and
thin and have small density; this is Mahler's criterion (Lecture~4) in its simplest form.

\begin{figure}[ht]
\centering
\begin{tikzpicture}[scale=2.3]
  \fill[addcol!12] (-0.5,3.0) -- (-0.5,0.866) arc (120:60:1) -- (0.5,3.0) -- cycle;
  \draw[->] (-1.7,0) -- (1.8,0) node[right]{$\operatorname{Re}\tau$};
  \draw[->] (0,-0.05) -- (0,3.15);
  \foreach \c in {-1,1} \draw[gray,dashed] (\c,0) ++(0:1) arc (0:180:1);
  \draw[gray,dashed] (1,0) arc (0:180:1);
  \foreach \c in {-1.5,1.5} \draw[gray,dashed] (\c,0) -- (\c,3.0);
  \draw[thick] (-0.5,3.0) -- (-0.5,0.866) arc (120:60:1) -- (0.5,3.0);
  \draw[flagcol,thin] (-1.6,0.866) -- (1.7,0.866) node[below right,font=\scriptsize]{$\dens=\frac{\pi}{\sqrt{12}}$};
  \draw[flagcol,thin] (-1.6,1) -- (1.7,1) node[above right,font=\scriptsize]{$\dens=\frac{\pi}{4}$};
  \draw[flagcol,thin] (-1.6,2) -- (1.7,2) node[right,font=\scriptsize]{$\dens=\frac{\pi}{8}$};
  \fill (0,1) circle (0.025) node[above left,font=\small]{$i$};
  \fill (0.5,0.866) circle (0.025) node[below right,font=\small]{$\rho$};
  \fill (-0.5,0.866) circle (0.025) node[below left,font=\small]{$\rho^2$};
  \foreach \x/\l in {-1/{-1},-0.5/{-\frac12},0.5/{\frac12},1/{1}} \node[below,font=\scriptsize] at (\x,0) {$\l$};
  \node[font=\small] at (0,2.6) {$\mathcal F$};
  \node[font=\scriptsize,gray,align=center] at (1.2,2.75) {long thin lattices\\(cusp)};
  % little lattice pictures
  \begin{scope}[shift={(-1.35,1.55)},scale=0.13]
    \foreach \i in {-2,...,2} \foreach \j in {-1,...,1} \fill ({\i+0.5*\j},{0.866*\j}) circle (0.18);
  \end{scope}
  \draw[gray,-{Stealth}] (-1.1,1.4) -- (0.42,0.9);
  \begin{scope}[shift={(1.25,1.55)},scale=0.13]
    \foreach \i in {-2,...,2} \foreach \j in {-1,...,1} \fill (\i,\j) circle (0.18);
  \end{scope}
  \draw[gray,-{Stealth}] (1.0,1.45) -- (0.06,1.03);
\end{tikzpicture}
\caption{The reduced domain $\mathcal F=\{\abs{\operatorname{Re}\tau}\le\frac12,\ \abs\tau\ge1\}$ for
binary forms, with some of its $\SL_2(\Z)$-translates dashed. Each point is a planar lattice with basis
$1,\tau$, up to rotation and scaling. The density $\pi/(4\operatorname{Im}\tau)$ is constant on
horizontal lines and is largest at the corners $\rho,\rho^2$, the hexagonal lattice.}
\label{NT:fig:modular}
\end{figure}

\begin{remark}[Class numbers]\label{NT:rem:classno}
For integral forms ($a,b,c\in\Z$) the discriminant $D=b^2-4ac<0$ is an invariant, and reduced forms
satisfy $3a^2\le -D$. So there are only finitely many reduced forms of a given discriminant. The number of
\emph{primitive} ones ($\gcd(a,b,c)=1$), with the usual boundary convention ($b\ge0$ if $\abs b=a$ or $a=c$),
is the \emph{class number} $h(D)$. For instance, $h(-4)=1$ (only $x^2+y^2$) and $h(-20)=2$ ($x^2+5y^2$ and $2x^2+2xy+3y^2$). Gauss
asked which $D$ have $h(D)=1$; see the Gallery.
\end{remark}

%======================================================================
\section{Counting lattice points: the Gauss circle problem}
\label{NT:sec:circle}

The density of a lattice packing is a count in disguise: by Exercise~\ref{L1:ex:periodic-density} of
Lecture~1, the number of lattice points in a large ball is
\[
  N_\Lat(r)=\#\{x\in\Lat:\abs x\le r\}\;\sim\;\frac{V_nr^n}{\covol(\Lat)},
  \qquad\text{so}\qquad
  \dens(\Lat)=\lim_{r\to\infty}\frac{N_\Lat(r)\,V_n(\lambda_1/2)^n}{V_nr^n}.
\]
How good is the approximation? For $\Lat=\Z^2$ this is \emph{Gauss's circle problem}. Write
\[
  N(r)=\#\{(x,y)\in\Z^2:x^2+y^2\le r^2\}=\pi r^2+E(r).
\]

\begin{proposition}[Gauss]\label{NT:prop:gauss}
$\abs{E(r)}\le\pi\big(\sqrt2\,r+\tfrac12\big)$ for $r\ge1$; in particular $E(r)=O(r)$.
\end{proposition}

\begin{proof}
Attach to each lattice point $p$ the unit square $p+[-\frac12,\frac12]^2$; these squares tile the plane
and each has diameter $\sqrt2$. The squares of the $N(r)$ points in the closed disk of radius $r$ lie
in the disk of radius $r+\frac{\sqrt2}2$. Conversely, every point of the disk of radius
$r-\frac{\sqrt2}2$ lies in a square whose centre is in the disk of radius $r$. Comparing areas,
\[
  \pi\big(r-\tfrac{\sqrt2}{2}\big)^2\le N(r)\le\pi\big(r+\tfrac{\sqrt2}{2}\big)^2,
\]
and $\pi(r\pm\frac{\sqrt2}2)^2-\pi r^2=\pm\pi\sqrt2\,r+\frac\pi2$.
\end{proof}

\begin{figure}[ht]
\centering
\begin{tikzpicture}[scale=0.42]
  \def\rr{4.3}
  \foreach \i in {-5,...,5} \foreach \j in {-5,...,5} {
    \pgfmathtruncatemacro{\inside}{(\i*\i+\j*\j<=18.49)?1:0}
    \ifnum\inside=1
      \fill[addcol!15] (\i-0.5,\j-0.5) rectangle (\i+0.5,\j+0.5);
      \draw[addcol!50,very thin] (\i-0.5,\j-0.5) rectangle (\i+0.5,\j+0.5);
      \fill (\i,\j) circle (2.6pt);
    \else
      \fill[gray!60] (\i,\j) circle (1.6pt);
    \fi
  }
  \draw[thick] (0,0) circle (\rr);
  \draw[flagcol,dashed] (0,0) circle ({\rr+0.7071});
  \draw[flagcol,dashed] (0,0) circle ({\rr-0.7071});
  \draw[-{Stealth}] (0,0) -- (30:\rr) node[pos=0.6,above,font=\small]{$r$};
  \node[font=\small,flagcol] at (6.9,3.6) {$r\pm\tfrac{\sqrt2}{2}$};
\end{tikzpicture}
\caption{Gauss's argument for $r=4.3$: the $N(r)=61$ unit squares centred at the lattice points in the
disk are trapped between the disks of radii $r\pm\frac{\sqrt2}{2}$, so $\abs{N(r)-\pi r^2}=O(r)$.}
\label{NT:fig:gauss}
\end{figure}

The same argument, with a fundamental parallelepiped in place of the square, gives
$N_\Lat(r)=V_nr^n/\covol(\Lat)+O(r^{n-1})$ for every lattice and every dimension, and also for any
ellipsoid $\{Q(x)\le r^2\}$ in place of the ball. So the error term is a ``boundary layer'', as in the
density remarks of Lecture~1. The question is how much the lattice points conspire along the boundary.

\paragraph{Representation numbers.} Let $r_2(m)=\#\{(x,y)\in\Z^2:x^2+y^2=m\}$. Then
$N(\sqrt X)=1+\sum_{1\le m\le X}r_2(m)$, so Proposition~\ref{NT:prop:gauss} says that \emph{the average
value of $r_2(m)$ is $\pi$}. At the same time $r_2(m)$ is given exactly by a formula of Jacobi (1829):
with $\chi(d)=0,1,-1$ for $d$ even, $d\equiv1$, $d\equiv3\pmod4$,
\begin{equation}\label{NT:eq:r2}
  r_2(m)=4\sum_{d\mid m}\chi(d)=4\big(d_1(m)-d_3(m)\big),
\end{equation}
where $d_i(m)$ counts the divisors of $m$ that are $\equiv i\pmod 4$. (For example $r_2(5)=8$ and
$r_2(3)=0$; it was checked numerically for all $m<2000$ in preparing these notes.) Combining the
geometric count with the arithmetic formula gives a lattice point proof of Leibniz's series.

\begin{proposition}[Lattice points prove Leibniz]\label{NT:prop:leibniz}
$1-\frac13+\frac15-\frac17+\dots=\frac\pi4$.
\end{proposition}

\begin{proof}
Put $L=\sum_{d\ge1}\chi(d)/d$ and $S(t)=\sum_{d\le t}\chi(d)\in\{0,1\}$. By \eqref{NT:eq:r2},
$N(\sqrt X)-1=4\sum_{dm\le X}\chi(d)$, a sum over the lattice points under the hyperbola $dm=X$. With
$y=\sqrt X$, Dirichlet's hyperbola method splits it as
\[
  \sum_{dm\le X}\chi(d)=\sum_{d\le y}\chi(d)\Big\lfloor\frac Xd\Big\rfloor+\sum_{m\le y}S\Big(\frac Xm\Big)-S(y)\lfloor y\rfloor .
\]
The last two terms are $O(y)$. In the first, $\lfloor X/d\rfloor=X/d+O(1)$, and the alternating series
gives $\sum_{d\le y}\chi(d)/d=L+O(1/y)$. Hence $N(\sqrt X)=4LX+O(\sqrt X)$. Proposition~\ref{NT:prop:gauss}
gives $N(\sqrt X)=\pi X+O(\sqrt X)$. Divide by $X$ and let $X\to\infty$.
\end{proof}

\paragraph{How big is $E(r)$?} Gauss's $O(r)$ uses nothing about circles except that they are curves
of finite length. Curvature helps: Sierpi\'nski (1906) proved $E(r)=O(r^{2/3})$, and the exponent has
been lowered slowly since. The conjecture, usually attributed to Hardy, is that
\[
  E(r)=O_\varepsilon\big(r^{1/2+\varepsilon}\big)\qquad\text{for every }\varepsilon>0,
\]
and this would be optimal: Hardy and, independently, Landau showed in 1915 that $E(r)$ is not
$o(r^{1/2})$, and Hardy (1916) that it is \emph{not} $o(r^{1/2}(\log r)^{1/4})$, so the
exponent $\frac12$ cannot be lowered. Sharper logarithmic factors are due to Hafner (1981) and
Soundararajan (2003). The exponent $\frac12$ is also the right one on average: Cram\'er (Math.\ Z.\ 15,
1922) showed
that $\frac{1}{R}\int_0^RE(r)^2\,dr$ grows like a constant times $R$, so a typical $\abs{E(r)}$ is of
size $r^{1/2}$. The current upper bounds, as exponents of $r$, are
\begin{center}
\renewcommand{\arraystretch}{1.15}
\begin{tabular}{lll}
\toprule
Gauss (c.~1800) & $1$ & elementary\\
Sierpi\'nski (1906) & $2/3\approx0.6667$ & \\
Huxley (2003) & $131/208\approx0.62981$ & Proc.\ London Math.\ Soc.\ 87\\
Bourgain--Watt (2017) & $517/824\approx0.62743$ & arXiv:1709.04340 (claimed)\\
Li--Yang (2023) & $\approx0.628966$ & arXiv:2308.14859\\
conjecture & $1/2+\varepsilon$ & \\
\bottomrule
\end{tabular}
\end{center}
Huxley's exponent is the best one published in a refereed journal: as of 2026 both later papers are
unpublished arXiv preprints. Both build on the Bombieri--Iwaniec method, and Bourgain--Watt add
Bourgain's decoupling theory. The Bourgain--Watt preprint claims the smaller exponent, but Li and Yang
describe Huxley's bound as the best known before their work and present theirs as its first improvement.
So the record is not settled, and should be quoted with care.

\begin{remark}[Why a packing person should care]
The circle problem is the cleanest model of a phenomenon that appears all through these notes: the bulk
term of a count is a volume, and hence a density, while the correction is a boundary term whose size
depends on fine arithmetic. In Lecture~2 we removed the boundary term by letting a box grow. Here it is
the main object of study.
\end{remark}

%======================================================================
\section{Blichfeldt again: Minkowski's convex body theorem}
\label{NT:sec:minkowski}

In Lecture~2 Blichfeldt's method compared a pointwise bound with an average: if
$\sum_if(\varphi_i^{-1}x)\le1$ everywhere, then its average, which is proportional to the density, is at
most $1$. Fifteen years earlier Blichfeldt had used the same comparison the other way round.

\begin{lemma}[Blichfeldt 1914]\label{NT:lem:blichfeldt}
Let $S\subset\R^n$ be measurable with $\vol(S)>\covol(\Lat)$. Then $S$ contains two distinct points
$s,s'$ with $s-s'\in\Lat$.
\end{lemma}

\begin{proof}
Consider the periodization $F(x)=\sum_{\lambda\in\Lat}\chi_S(x+\lambda)$, a $\Lat$-periodic function.
Integrating over a fundamental domain $P$,
\[
  \int_PF(x)\,dx=\sum_{\lambda\in\Lat}\vol\big(S\cap(P+\lambda)\big)=\vol(S)>\vol(P).
\]
So $F$ cannot be $\le1$ everywhere: there is an $x$ with $F(x)\ge2$, that is, two different
$\lambda,\lambda'\in\Lat$ with $s=x+\lambda$ and $s'=x+\lambda'$ in $S$. Then $s-s'=\lambda-\lambda'\in\Lat\setminus0$.
\end{proof}

\begin{figure}[ht]
\centering
\begin{tikzpicture}[scale=1.25]
  \def\blob{plot[smooth cycle,tension=0.8] coordinates {(0.25,0.35) (1.3,0.15) (2.35,0.55) (2.6,1.45) (1.75,2.05) (0.8,1.8) (0.15,1.2)}}
  \foreach \i in {0,...,3} \draw[gray!60] (\i,0) -- (\i,2.3);
  \foreach \j in {0,...,2} \draw[gray!60] (-0.1,\j) -- (3.1,\j);
  \foreach \i in {0,...,3} \foreach \j in {0,...,2} \fill[gray] (\i,\j) circle (0.025);
  \fill[addcol!25] \blob; \draw[addcol,thick] \blob;
  \node[addcol,font=\small] at (1.4,1.05) {$S$};
  \draw[-{Stealth},thick] (3.35,1) -- (4.15,1) node[midway,above,font=\scriptsize]{mod $\Lat$};
  \begin{scope}[xshift=4.5cm,yshift=0cm,scale=2]
    \foreach \i in {0,1,2} \foreach \j in {0,1} {
      \begin{scope}[shift={(-\i,-\j)}]
        \clip (\i,\j) rectangle (\i+1,\j+1);
        \fill[addcol,opacity=0.28] \blob;
        \draw[addcol] \blob;
      \end{scope}}
    \draw[thick] (0,0) rectangle (1,1);
    \node[below,font=\small] at (0.5,0) {$P=[0,1)^2$};
  \end{scope}
\end{tikzpicture}
\caption{Blichfeldt's lemma for $\Lat=\Z^2$. The pieces of $S$ in the different unit squares are
translated into one fundamental square $P$. Their total area is $\vol(S)>1=\vol(P)$, so two pieces
overlap (the darker regions). A point of overlap comes from two points of $S$ that differ by a lattice
vector.}
\label{NT:fig:blichfeldt}
\end{figure}

Notice the parallel. In Lecture~2 the periodized function was bounded by $1$ (packing condition), and
the conclusion was that its average, and so the density, is at most $1$. Here the average exceeds $1$, so
the function is somewhere at least $2$: the translates $S+\lambda$ are \emph{not} a packing. Both are
pigeonhole arguments for the periodization of a function.

\begin{theorem}[Minkowski's convex body theorem, 1889/1896]\label{NT:thm:minkowski}
Let $K\subset\R^n$ be convex and symmetric ($K=-K$). If $\vol(K)>2^n\covol(\Lat)$, then $K$ contains a
nonzero point of $\Lat$. If $K$ is also compact, $\vol(K)\ge2^n\covol(\Lat)$ suffices.
\end{theorem}

\begin{proof}
Apply Lemma~\ref{NT:lem:blichfeldt} to $S=\frac12K$, of volume $2^{-n}\vol(K)>\covol(\Lat)$. This
gives $s\ne s'$ in $\frac12K$ with $s-s'\in\Lat$. By symmetry $-s'\in\frac12K$, and by convexity
$\frac12(2s)+\frac12(-2s')=s-s'\in K$. For compact $K$ apply the first part to $(1+\varepsilon)K$ for
each $\varepsilon>0$. Each of these contains a nonzero lattice point, and they all lie in the bounded set
$2K$, which contains only finitely many lattice points. So one nonzero point lies in $(1+\varepsilon)K$
for arbitrarily small $\varepsilon$, hence in $K$.
\end{proof}

\begin{remark}[Minkowski and the trivial bound]
Applied to the ball $K=\Ball^n(\lambda_1)$, Minkowski gives $V_n\lambda_1^n\le2^n\covol(\Lat)$, which
is $\dens(\Lat)\le1$, the trivial bound that the balls of radius $\lambda_1/2$ are disjoint. (So the two
estimates quoted in Lecture~4 are actually the same inequality.) For sphere packing Minkowski's theorem
is therefore weak. Its power lies in two other directions: in arithmetic, where the lattice is
cut out by congruences, and in the bodies $K$ other than balls. Here are the two classic examples.
\end{remark}

\begin{theorem}[Fermat]\label{NT:thm:fermat}
Every prime $p\equiv1\pmod4$ is a sum of two squares.
\end{theorem}

\begin{proof}
Since $p\equiv1\pmod4$, $-1$ is a square modulo $p$; for instance $u=\big(\frac{p-1}2\big)!$ satisfies
$u^2\equiv-1$ by Wilson's theorem. Let
\[
  \Lat=\{(x,y)\in\Z^2:\ y\equiv ux\pmod p\},
\]
a sublattice of index $p$ in $\Z^2$, so $\covol(\Lat)=p$. The open disk of radius $\sqrt{2p}$ has area
$2\pi p>4p$, so by Theorem~\ref{NT:thm:minkowski} it contains some $(x,y)\in\Lat\setminus0$. Then
$0<x^2+y^2<2p$ and $x^2+y^2\equiv x^2(1+u^2)\equiv0\pmod p$. Hence $x^2+y^2=p$.
\end{proof}

\begin{figure}[ht]
\centering
\begin{tikzpicture}[scale=0.4]
  \foreach \x in {-6,...,6} \foreach \y in {-6,...,6} \fill[gray!40] (\x,\y) circle (1.5pt);
  \foreach \x in {-6,...,6} \foreach \y in {-6,...,6} {
    \pgfmathtruncatemacro{\t}{mod(\y-5*\x+130,13)}
    \ifnum\t=0 \fill (\x,\y) circle (4pt); \fi }
  \draw[thick,addcol] (0,0) circle ({sqrt(26)});
  \draw[-{Stealth},flagcol,thick] (0,0) -- (2,-3) node[below right,font=\scriptsize]{$(2,-3)$};
  \draw[-{Stealth},flagcol,thick] (0,0) -- (3,2) node[right,font=\scriptsize]{$(3,2)$};
  \node[addcol,font=\small] at (5.6,4.9) {$\abs v=\sqrt{26}$};
\end{tikzpicture}
\caption{Fermat by Minkowski for $p=13$, $u=5$ ($5^2\equiv-1$). The lattice
$\{y\equiv5x\bmod13\}$ (black) has covolume $13$. The disk of area $26\pi>4\cdot13$ must contain
lattice vectors, and they have $x^2+y^2=13$. Here the lattice is a square lattice of side $\sqrt{13}$,
with reduced basis $(3,2)$, $(2,-3)$.}
\label{NT:fig:fermat}
\end{figure}

\begin{theorem}[Lagrange's four square theorem, 1770]\label{NT:thm:lagrange4}
Every positive integer is a sum of four squares.
\end{theorem}

\begin{proof}
By Euler's identity, which expresses $(a_1^2+\dots+a_4^2)(b_1^2+\dots+b_4^2)$ as a sum of four squares
(it is the multiplicativity of the quaternion norm), it suffices to treat primes, and $2=1^2+1^2$. Let
$p$ be an odd prime. The sets $\{a^2\bmod p\}$ and $\{-1-b^2\bmod p\}$ each have $\frac{p+1}2$
elements, so they meet: $a^2+b^2+1\equiv0\pmod p$ for some $a,b$. Let
\[
  \Lat=\{x\in\Z^4:\ x_3\equiv ax_1+bx_2,\ \ x_4\equiv bx_1-ax_2\pmod p\},
\]
of index $p^2$ in $\Z^4$. For $x\in\Lat$,
\[
  x_1^2+x_2^2+x_3^2+x_4^2\equiv(x_1^2+x_2^2)(1+a^2+b^2)\equiv0\pmod p .
\]
The open ball of radius $\sqrt{2p}$ in $\R^4$ has volume $\frac{\pi^2}{2}(2p)^2=2\pi^2p^2>16p^2$, so
it contains a nonzero $x\in\Lat$, and then $0<\abs x^2<2p$ and $p\mid\abs x^2$ force $\abs x^2=p$.
\end{proof}

Three squares are more delicate. There is no product identity for sums of three squares, and the known
geometric proof (Ankeny, 1957) needs Dirichlet's theorem on primes in arithmetic progressions in
addition to Minkowski's theorem.

\begin{theorem}[Legendre 1798; Gauss 1801]\label{NT:thm:legendre}
A positive integer is a sum of three squares if and only if it is not of the form $4^a(8b+7)$.
\end{theorem}

The ``only if'' half is a computation mod $8$. The ``if'' half needs the theory of ternary forms (or
Dirichlet's theorem on primes in progressions). Which other forms represent \emph{every} positive integer?
Ramanujan (1917) listed $55$ diagonal quaternary forms $ax^2+by^2+cz^2+dw^2$ with this property. One of
them, $x^2+2y^2+5z^2+5w^2$, fails to represent $15$, and the other $54$ are exactly the universal ones. The definitive answer is surprisingly finite; see the $15$ and $290$ theorems in
the Gallery.

%======================================================================
\section{Theta series and modular forms}
\label{NT:sec:theta}

All the representation numbers of a lattice can be packaged into one function.

\begin{definition}\label{NT:def:theta}
The \emph{theta series} of $\Lat$ is
\[
  \Theta_\Lat(\tau)=\sum_{x\in\Lat}e^{\pi i\tau\abs x^2}=\sum_{t\ge0}N_t\,e^{\pi i\tau t},
  \qquad \tau\in\mathbb H,
\]
where $N_t$ is the number of vectors of squared length $t$. The series converges absolutely and locally
uniformly on $\mathbb H$, because $\abs{e^{\pi i\tau\abs x^2}}=e^{-\pi\operatorname{Im}\tau\abs x^2}$.
\end{definition}

For example $\Theta_{\Z}(\tau)=\sum_{k\in\Z}q^{k^2}$ with $q=e^{\pi i\tau}$ is Jacobi's theta function,
and $\Theta_{\Lat\oplus\Lat'}=\Theta_\Lat\Theta_{\Lat'}$, so
$\Theta_{\Z^2}=\Theta_\Z^2=\sum_mr_2(m)q^m$. Formula \eqref{NT:eq:r2} and Jacobi's four square formula
\begin{equation}\label{NT:eq:r4}
  r_4(m)=8\sum_{d\mid m,\ 4\nmid d}d
\end{equation}
are identities between $q$-series (both were checked numerically for small $m$). The tool that makes
theta series powerful is the following.

\begin{theorem}[Poisson summation over a lattice]\label{NT:thm:poisson}
Let $f\colon\R^n\to\C$ be continuous and satisfy $\abs{f(x)},\abs{\hat f(x)}\le C(1+\abs x)^{-n-\eta}$ for some $C,\eta>0$,
where $\hat f(y)=\int_{\R^n}f(x)e^{-2\pi i\ip xy}dx$. Then for every $v\in\R^n$,
\[
  \sum_{x\in\Lat}f(x+v)=\frac{1}{\covol(\Lat)}\sum_{y\in\Lat^*}\hat f(y)\,e^{2\pi i\ip yv},
  \qquad \Lat^*=\{y\in\R^n:\ip xy\in\Z\ \text{for all }x\in\Lat\}.
\]
Both sides converge absolutely.
\end{theorem}

\begin{proof}[Sketch]
The left side $F(v)$ is a continuous $\Lat$-periodic function, so it is a function on the torus
$\R^n/\Lat$, whose characters are $v\mapsto e^{2\pi i\ip yv}$, $y\in\Lat^*$. Its Fourier coefficients are
\[
  \frac{1}{\covol(\Lat)}\int_{\R^n/\Lat}F(v)e^{-2\pi i\ip yv}dv
  =\frac{1}{\covol(\Lat)}\int_{\R^n}f(v)e^{-2\pi i\ip yv}dv=\frac{\hat f(y)}{\covol(\Lat)},
\]
by unfolding the sum against the integral. These coefficients are absolutely summable by the decay of
$\hat f$, so the Fourier series converges uniformly to $F$.
\end{proof}

The \emph{dual lattice} $\Lat^*$ has $\covol(\Lat^*)=1/\covol(\Lat)$ and Gram matrix $G^{-1}$. $\Lat$ is
integral iff $\Lat\subseteq\Lat^*$, and \emph{unimodular} iff $\Lat=\Lat^*$ (equivalently, it is integral with
$\covol=1$). It is \emph{even} if all $\abs x^2$ are even. $\Z^n$ is unimodular but odd; $E_8$ is even
unimodular: it is integral and has determinant $1$ (Lecture~6).

Apply Theorem~\ref{NT:thm:poisson} to a Gaussian. For $\tau\in\mathbb H$ the function
$f(x)=e^{\pi i\tau\abs x^2}$ has $\hat f(y)=(\tau/i)^{-n/2}e^{\pi i(-1/\tau)\abs y^2}$ (principal branch;
for $\tau=it$ this is the familiar fact that $e^{-\pi t\abs x^2}$ has Fourier transform
$t^{-n/2}e^{-\pi\abs y^2/t}$, and both sides are holomorphic in $\tau$). With $v=0$ this gives $\Theta_\Lat(\tau)=\covol(\Lat)^{-1}(i/\tau)^{n/2}\,\Theta_{\Lat^*}(-1/\tau)$; now replace $\tau$ by $-1/\tau$.

\begin{corollary}[Theta transformation formula]\label{NT:cor:thetatrans}
$\displaystyle \Theta_\Lat\Big(-\frac1\tau\Big)=\Big(\frac\tau i\Big)^{n/2}\frac{1}{\covol(\Lat)}\,\Theta_{\Lat^*}(\tau).$
\end{corollary}

So the theta series of a lattice controls that of its dual, through the substitution $\tau\mapsto-1/\tau$.
For an even unimodular lattice we get two symmetries:
\begin{equation}\label{NT:eq:modular}
  \Theta_\Lat(\tau+1)=\Theta_\Lat(\tau),\qquad
  \Theta_\Lat\Big(-\frac1\tau\Big)=\Big(\frac\tau i\Big)^{n/2}\Theta_\Lat(\tau).
\end{equation}
The maps $T\colon\tau\mapsto\tau+1$ and $S\colon\tau\mapsto-1/\tau$ generate $\SL_2(\Z)$, the group
whose fundamental domain we drew in Figure~\ref{NT:fig:modular}. The fact that \emph{lattices in the
plane} are parametrized by $\SL_2(\Z)\backslash\mathbb H$ and the fact that \emph{theta series of lattices
in any dimension} transform under $\SL_2(\Z)$ are two separate appearances of the same group.

\begin{theorem}\label{NT:thm:8n}
If $\R^n$ contains an even unimodular lattice, then $8\mid n$.
\end{theorem}

\begin{proof}
Let $f=\Theta_\Lat$. Since $\Theta_{\Lat\oplus\Lat}=f^2$ and $\Lat\oplus\Lat$ is again even unimodular,
we may pass to $\Lat\oplus\Lat$ or to $\Lat^{\oplus4}$. So if $8\nmid n$ we may assume $n\equiv4\pmod 8$.
Put $k=n/2\equiv2\pmod4$. Then $(\tau/i)^{k}=i^{-k}\tau^k=-\tau^k$, and \eqref{NT:eq:modular} reads
$f(\tau+1)=f(\tau)$, $f(-1/\tau)=-\tau^kf(\tau)$. For $\gamma=\left(\begin{smallmatrix}a&b\\c&d\end{smallmatrix}\right)\in\SL_2(\Z)$
define $(f|\gamma)(\tau)=(c\tau+d)^{-k}f(\gamma\tau)$. The chain rule for $\gamma\tau$ shows
that $(f|\gamma)|\gamma'=f|(\gamma\gamma')$, and $f|(-I)=f$ since $k$ is even. With
$T=\left(\begin{smallmatrix}1&1\\0&1\end{smallmatrix}\right)$ and $S=\left(\begin{smallmatrix}0&-1\\1&0\end{smallmatrix}\right)$
we have $f|T=f$ and $f|S=-f$. A direct multiplication gives $(ST)^3=-I$. Therefore
\[
  f=f|(-I)=f|STSTST=(-1)^3f=-f,
\]
so $f\equiv0$. This contradicts $f(\tau)\to1$ as $\operatorname{Im}\tau\to\infty$ (only $x=0$ survives).
\end{proof}

Conversely, even unimodular lattices exist in every dimension $8k$: $E_8^{\oplus k}$, or
$D_{8k}^+=D_{8k}\cup(D_{8k}+\frac12\mathbf 1)$ (Exercise~\ref{NT:ex:D16}; $E_8=D_8^+$ by Lecture~6).

\subsection*{The theta series of \texorpdfstring{$E_8$}{E8}}

For $8\mid n$ the factor $(\tau/i)^{n/2}$ equals $\tau^{n/2}$, and \eqref{NT:eq:modular} says that
$\Theta_\Lat$ is a \emph{modular form of weight $n/2$} for $\SL_2(\Z)$: a holomorphic function on
$\mathbb H$, bounded as $\operatorname{Im}\tau\to\infty$, with $f(\gamma\tau)=(c\tau+d)^{n/2}f(\tau)$.
Write $\mathsf q=e^{2\pi i\tau}$ (so $\mathsf q=q^2$). The basic theorem on such forms (see Serre,
\emph{A Course in Arithmetic}, Ch.~VII) is that the space $M_k$ of modular forms of weight $k$ is
finite-dimensional, with
\[
  M_4=\C E_4,\qquad M_8=\C E_4^2,\qquad M_{12}=\C E_4^3\oplus\C\Delta,
\]
where $E_4=1+240\sum_{m\ge1}\sigma_3(m)\mathsf q^m$ is the Eisenstein series
($\sigma_3(m)=\sum_{d\mid m}d^3$) and $\Delta=\mathsf q\prod_{m\ge1}(1-\mathsf q^m)^{24}$. Since
$\Theta_{E_8}=\sum_mN_{2m}\mathsf q^m$ lies in $M_4$ and has constant term $1$,
\begin{equation}\label{NT:eq:E8theta}
  \Theta_{E_8}=E_4,\qquad\text{i.e.}\qquad
  \#\{x\in E_8:\abs x^2=2m\}=240\,\sigma_3(m).
\end{equation}
For $m=1$ this is $240$, the root count from Lecture~6 and the dimension~$8$ kissing number. For $m=2$
it predicts $240\cdot9=2160$ vectors of norm $4$, and for $m=3$, $240\cdot28=6720$. A brute-force count
of $E_8$ vectors confirms $240,\,2160,\,6720$. \emph{A single modular form knows every shell of $E_8$.}

\subsection*{Hearing the shape of a drum, and Niemeier's lattices}

The flat torus $\R^n/\Lat$ has Laplace eigenfunctions $e^{2\pi i\ip yx}$, $y\in\Lat^*$, with eigenvalues
$4\pi^2\abs y^2$. So its spectrum, with multiplicities, is recorded by $\Theta_{\Lat^*}$.

\begin{theorem}[Milnor 1964]\label{NT:thm:milnor}
The $16$-dimensional flat tori $\R^{16}/(E_8\oplus E_8)$ and $\R^{16}/D_{16}^+$ have the same Laplace
spectrum but are not isometric.
\end{theorem}

\begin{proof}
Both lattices are even unimodular, so both theta series lie in $M_8=\C E_4^2$ and have constant term $1$.
So they are equal, and the lattices are self-dual, which makes the spectra equal. (Numerically, both have
$480$ vectors of norm $2$ and $61920$ of norm $4$.) An isometry of the tori would lift to an isometry of
$\R^{16}$ carrying one lattice onto the other. But the $480$ norm-$2$ vectors of $E_8\oplus E_8$ span it,
while those of $D_{16}^+$ (the vectors $\pm e_i\pm e_j$) span only the index-$2$ sublattice $D_{16}$.
\end{proof}

This was the first counterexample to the question ``can one hear the shape of a drum?'', two years
before Kac named the question. In dimension $24$, $M_{12}$ is two-dimensional, so $\Theta_\Lat$ is
determined by its constant term and by the number of roots. Taking zero roots gives
\[
  \Theta_{\Lambda_{24}}=E_4^3-720\Delta=1+196560\,\mathsf q^2+16773120\,\mathsf q^3+\cdots,
\]
and $196560$ is the kissing number of the Leech lattice. Niemeier (1973) classified the even unimodular
lattices in dimension $24$. There are exactly $24$ of them, determined by their root systems (each one of
rank $24$ or empty), and the Leech lattice is the unique one without roots (Conway 1969).

%======================================================================
\section{Poisson summation as the engine: LP bounds}
\label{NT:sec:LP}

Lecture~3 mentioned that the solutions in dimensions $8$ and $24$ use ``linear programming duality'' and a
``magic'' auxiliary function. With Poisson summation in hand, the argument for lattices takes only a few lines.

\begin{definition}
Call $f\colon\R^n\to\R$ \emph{admissible} if $f$ and $\hat f$ satisfy the decay hypothesis of
Theorem~\ref{NT:thm:poisson}, and $f$ is even (so $\hat f$ is real).
\end{definition}

\begin{theorem}[Cohn--Elkies bound for lattices]\label{NT:thm:CElattice}
Let $f$ be admissible with $f(0)=\hat f(0)>0$, and suppose
\begin{enumerate}[label=(\roman*),nosep]
\item $f(x)\le0$ for $\abs x\ge r$, and
\item $\hat f(y)\ge0$ for all $y$.
\end{enumerate}
Then every lattice packing in $\R^n$ has density at most $V_n(r/2)^n$.
\end{theorem}

\begin{proof}
Scale the lattice so that $\lambda_1(\Lat)=r$. Apply Poisson summation with $v=0$:
\[
  f(0)\;\ge\;f(0)+\sum_{x\in\Lat\setminus0}f(x)\;=\;\frac{1}{\covol(\Lat)}\sum_{y\in\Lat^*}\hat f(y)
  \;\ge\;\frac{\hat f(0)}{\covol(\Lat)}.
\]
The first inequality is (i), since every $x\ne0$ has $\abs x\ge r$, and the last is (ii), dropping
all $y\ne0$. Hence $\covol(\Lat)\ge\hat f(0)/f(0)=1$, and $\dens(\Lat)=V_n(r/2)^n/\covol(\Lat)\le V_n(r/2)^n$.
\end{proof}

The full theorem drops the lattice assumption at almost no cost.

\begin{theorem}[Cohn--Elkies 2003]\label{NT:thm:CE}
Under the same hypotheses, every sphere packing in $\R^n$ has density at most $V_n(r/2)^n$.
\end{theorem}

\begin{proof}
By Exercise~\ref{L1:ex:periodic-approx} of Lecture~1 it suffices to treat periodic packings, with
centres $\bigcup_{j=1}^N(\Lat+v_j)$, the $v_j$ pairwise incongruent modulo $\Lat$, scaled to have
minimal distance $r$. Sum Poisson summation over all pairs
$(j,k)$ with $v=v_j-v_k$:
\[
  \sum_{j,k=1}^N\sum_{x\in\Lat}f(x+v_j-v_k)
  =\frac{1}{\covol(\Lat)}\sum_{y\in\Lat^*}\hat f(y)\,\Big|\sum_{j=1}^Ne^{2\pi i\ip y{v_j}}\Big|^2 .
\]
On the left, $x+v_j-v_k$ is the difference of two centres. It is $0$ exactly when $j=k$ and $x=0$, and
otherwise it has length at least $r$. So the left side is at most $Nf(0)$. On the right every term is
$\ge0$, and the $y=0$ term is $N^2\hat f(0)/\covol(\Lat)$. So $N/\covol(\Lat)\le f(0)/\hat f(0)=1$. The
density is $(N/\covol(\Lat))V_n(r/2)^n$.
\end{proof}

\begin{remark}[Blichfeldt's gauge versus Cohn--Elkies]\label{NT:rem:blichfeldt-vs-CE}
Recall the proof that $f_0=1-\frac12\abs x^2$ is a gauge (Lemma~\ref{L2:lem:f0}). It rested on the identity
\[
  \sum_{i<j}\abs{x_i-x_j}^2=m\sum_i\abs{x_i}^2-\Big|\sum_ix_i\Big|^2,
\]
and the inequality came from \emph{throwing away a square}, $\abs{\sum_ix_i}^2\ge0$. The proof of
Theorem~\ref{NT:thm:CE} throws away the squares $\hat f(y)\abs{\sum_je^{2\pi i\ip y{v_j}}}^2\ge0$, one for
each frequency $y\ne0$. Blichfeldt uses one quadratic ``moment'', while Cohn--Elkies use all the
exponentials. There is also a difference in what the auxiliary function must satisfy:
\begin{itemize}[nosep]
\item a Blichfeldt gauge is a condition in physical space only: $f\ge0$ and
  $\sum_if(\varphi_i^{-1}x)\le1$ for every packing;
\item a Cohn--Elkies function has a sign condition in \emph{both} spaces: $f\le0$ beyond $r$, and
  $\hat f\ge0$ everywhere.
\end{itemize}
The uncertainty principle makes the second kind of condition hard to satisfy, and this is why the
optimal functions are rare and remarkable.
\end{remark}

\paragraph{When is the bound sharp?} Suppose $\Lat$ is unimodular, $\lambda_1(\Lat)=r$, and $f$ gives
equality. Then both inequalities in the proof of Theorem~\ref{NT:thm:CElattice} are equalities, so
\[
  f(x)=0\ \text{ for all }x\in\Lat\setminus0,\qquad \hat f(y)=0\ \text{ for all }y\in\Lat^*\setminus0.
\]
For $E_8$ ($r=\sqrt2$) this says that $f$ and $\hat f$ vanish on all spheres $\abs x=\sqrt{2m}$, $m\ge1$. Since
$f\le0$ beyond $\sqrt2$ and $\hat f\ge0$ everywhere, these zeros must have even order (they are double roots),
except that $f$ may change sign, with a simple root, at $\sqrt2$ itself.
Cohn and Elkies found numerically that such functions seem to exist in dimensions $8$ and $24$ (and
$2$), but could not construct them.

\paragraph{Viazovska's idea.} Look for $f$ radial and built out of Gaussians:
\[
  f(x)=\int_{\Gamma}\phi(z)\,e^{\pi iz\abs x^2}\,dz
\]
along a path $\Gamma$ in $\mathbb H$. In dimension $8$ the Fourier transform of $e^{\pi iz\abs x^2}$ is
$z^{-4}e^{\pi i(-1/z)\abs y^2}$ (the computation before Corollary~\ref{NT:cor:thetatrans}). Substituting
$w=-1/z$ shows that $\hat f$ is given by the same kind of integral with weight $\phi(-1/w)\,w^{2}$ in place of
$\phi(w)$, over the image path. So the Fourier transform acts on the \emph{weight} $\phi$ through
$z\mapsto-1/z$, and $f$ is a Fourier eigenfunction when $\phi$ transforms like a modular form of weight $-2$
(up to the change of path). Viazovska (2016) built two eigenfunctions of the Fourier transform, with
eigenvalues $\pm1$: one from a weakly holomorphic \emph{quasimodular} form for $\SL_2(\Z)$ (a rational
expression in $E_2,E_4,E_6$ and $\Delta$), and one from a weakly holomorphic modular form for the subgroup
$\Gamma(2)$ (built from Jacobi theta functions). A factor $\sin^2(\pi\abs x^2/2)$ supplies the
double roots at $\abs x^2\in2\Z$. A suitable linear combination of the two is the magic function, and it proves that $E_8$ is the densest
packing in $\R^8$. A week later, Cohn, Kumar, Miller, Radchenko and Viazovska did the same in dimension
$24$ for the Leech lattice. The theta series of Section~\ref{NT:sec:theta} and the magic functions use
the same symmetry $\tau\mapsto-1/\tau$, once applied to a lattice and once to an auxiliary function.

\paragraph{Universal optimality.} The same circle of ideas gives a much stronger statement (Gallery,
Theorem~\ref{NT:gal:universal}): $E_8$ and Leech minimize energy for every completely monotone function of
squared distance, Gaussians $e^{-\alpha\abs x^2}$ included. For Gaussian potentials, Poisson summation
turns the energy of a lattice into the energy of its dual. The proof needs a \emph{Fourier interpolation
formula}: a radial Schwartz function on $\R^8$ is determined by the values of $f$ and $\hat f$ (and
their radial derivatives) on the spheres $\abs x=\sqrt{2m}$.

%======================================================================
\section{Averaging over all lattices: Siegel and Minkowski--Hlawka}
\label{NT:sec:siegel}

So far we have looked at special lattices. At the opposite extreme, take a \emph{random} lattice. Let
$X_n=\SL_n(\R)/\SL_n(\Z)$ be the space of lattices of covolume $1$ in $\R^n$. It carries a unique
$\SL_n(\R)$-invariant probability measure $\mu$ (Minkowski, Siegel). For $n=2$, after quotienting by
rotations, $X_2$ becomes the modular surface of Figure~\ref{NT:fig:modular}, and $\mu$ becomes the
hyperbolic area $\frac{3}{\pi}\frac{dx\,dy}{y^2}$ restricted to $\mathcal F$.

\begin{theorem}[Siegel's mean value theorem, 1945]\label{NT:thm:siegel}
For $n\ge2$ and every integrable $f\colon\R^n\to\R$,
\[
  \int_{X_n}\sum_{x\in\Lat\setminus0}f(x)\,d\mu(\Lat)=\int_{\R^n}f(x)\,dx,
  \qquad
  \int_{X_n}\sum_{x\in\Lat\ \mathrm{primitive}}f(x)\,d\mu(\Lat)=\frac{1}{\zeta(n)}\int_{\R^n}f(x)\,dx .
\]
\end{theorem}

The two formulas are equivalent, because every nonzero vector is uniquely $kx$ with $k\ge1$ and $x$
primitive, and $\sum_kk^{-n}=\zeta(n)$. In words: \emph{on average, a lattice looks like Lebesgue measure.}
From this one gets a lower bound in a few lines.

\begin{corollary}[Minkowski--Hlawka]\label{NT:cor:MH}
For $n\ge2$, $\sup_\Lat\dens(\Lat)\ge\dfrac{2\zeta(n)}{2^n}$.
\end{corollary}

\begin{proof}
Let $B$ be a ball centred at $0$ with $\vol(B)<2\zeta(n)$, and apply the primitive form of
Theorem~\ref{NT:thm:siegel} to $f=\chi_B$. The average number of primitive vectors in $B$ is
$\vol(B)/\zeta(n)<2$. Primitive vectors come in pairs $\pm x$, so some $\Lat\in X_n$ has no primitive
vector in $B$, and hence no nonzero vector at all (if $kx\in B$ with $k\ge1$, then $x\in B$). Balls of
half the radius centred at $\Lat$ form a packing of density $\vol(B)/2^n$. Let
$\vol(B)\uparrow2\zeta(n)$.
\end{proof}

This is the ``$2\zeta(d)2^{-d}$'' of the remark after Proposition~\ref{L1:prop:saturated-bound} in
Lecture~1, and it is still within a polynomial factor of the best known lower bounds (for lattices,
Klartag's 2025 bound $c\,n^2\,2^{-n}$). Upper bounds are exponentially larger
(Kabatiansky--Levenshtein, Lecture~3). It is still unknown whether the best lattices in high dimension are
exponentially better than random ones.

Siegel's theorem treats $X_n$ as a probability space. Interlude~C treats it as a dynamical system: the
diagonal flow on $X_n$, Mahler's criterion as a statement about the cusp, and equidistribution results
that sharpen Siegel's average.

%======================================================================
\section{A dictionary}
\label{NT:sec:dictionary}

\begin{center}
\renewcommand{\arraystretch}{1.25}
\begin{tabular}{p{0.42\textwidth}p{0.50\textwidth}}
\toprule
Packings and lattices & Number theory\\
\midrule
lattice up to rotation & positive definite form up to $\GL_n(\Z)$\\
$\lambda_1^2$, $\covol^2$ & $\min Q$, $\det G$\\
densest lattice packing & Hermite's constant $\gamma_n$\\
planar lattices, Figure~\ref{NT:fig:modular} & $\SL_2(\Z)\backslash\mathbb H$, reduced binary forms\\
number of points in a ball & $\sum_{m\le r^2}r_Q(m)$; Gauss circle problem\\
kissing number of $\Lat$ & number of representations of $\min Q$\\
Blichfeldt's gauge (Lecture 2) & Blichfeldt's lemma, Minkowski's theorem\\
shells of $E_8$, Leech & coefficients of $E_4$, $E_4^3-720\Delta$\\
dual lattice $\Lat^*$ & $G^{-1}$; $\tau\mapsto-1/\tau$\\
LP bound, magic function & Poisson summation, (quasi)modular forms\\
random lattice & Siegel's mean value theorem\\
\bottomrule
\end{tabular}
\end{center}

%======================================================================
\section{Gallery of theorems}
\label{NT:sec:gallery}

Classical results first, then recent ones. Each comes with a sentence on why it belongs in this course.

\begin{theorem}[Hermite's inequality; Hermite 1850]\label{NT:gal:hermite}
$\gamma_n\le(4/3)^{(n-1)/2}$ for all $n\ge1$.
\end{theorem}
Published in Hermite's letters to Jacobi, J.\ reine angew.\ Math.\ 40 (1850). For $n=2$ it is
Theorem~\ref{NT:thm:gamma2}. The proof (Exercise~\ref{NT:ex:hermite}) is the row-stacking argument of
Lecture~1, run by induction on the dimension.

\begin{theorem}[The known Hermite constants]\label{NT:gal:hermiteknown}
The values of $\gamma_n$ in the table of Section~\ref{NT:sec:forms} are correct, and in each of these
dimensions the optimal lattice is unique up to similarity.
\end{theorem}
Lagrange (1773) and Gauss (1831, in his review of Seeber) for $n=2,3$; Korkine and Zolotarev, Math.\
Ann.\ 6 (1873) and 11 (1877), for $n=4,5$; Blichfeldt, Math.\ Z.\ 39 (1935), for $n=6,7,8$ (uniqueness
there is due to Vet\v{c}inkin, 1980); H.~Cohn and A.~Kumar, \emph{Optimality and uniqueness of the Leech
lattice among lattices}, Ann.\ of Math.\ 170 (2009), for $n=24$. The case $n=24$ uses the Cohn--Elkies
bound of Section~\ref{NT:sec:LP} together with an argument special to lattices. These are exactly the
lattices of Lectures~4--6.

\begin{theorem}[Voronoi 1908]\label{NT:gal:voronoi}
A positive definite form $Q$ is a local maximum of $\gamma(Q)$ (an \emph{extreme} form) if and only if it is
\emph{perfect} (it is the only form with its minimum value and its set of minimal vectors; equivalently,
the matrices $xx^T$, $x$ minimal, span the space of symmetric matrices) and \emph{eutactic} ($G^{-1}$ is a
combination of the matrices $xx^T$, $x$ minimal, with strictly positive coefficients). Up to equivalence and scaling there are finitely many perfect forms in each dimension.
\end{theorem}
G.~Voronoi, J.\ reine angew.\ Math.\ 133 (1908). The numbers of perfect forms in dimensions $2,\dots,8$
are $1,1,2,3,7,33,10916$ (dimension $7$: Jaquet-Chiffelle, Ann.\ Inst.\ Fourier 43 (1993); dimension $8$:
Dutour Sikiri\'c, Schürmann and Vallentin, Electron.\ Res.\ Announc.\ AMS 13 (2007)). This is a finite
algorithm for the lattice packing problem in each fixed dimension. Eutaxy is the lattice analogue of the
force balanced contact graphs of Lecture~10: the minimal vectors are the contacts, and the positive coefficients are
the stresses.

\begin{theorem}[Class number one; Heegner 1952, Baker 1966, Stark 1967]\label{NT:gal:classno}
The imaginary quadratic fields $\Q(\sqrt{-d})$ with class number $1$ are exactly those with
$d=1,2,3,7,11,19,43,67,163$.
\end{theorem}
Gauss conjectured the list. Heegner's proof (Math.\ Z.\ 56, 1952) was not accepted at first. Baker
(Mathematika 13, 1966) and Stark (Michigan Math.\ J.\ 14, 1967) gave independent proofs, and Heegner's
argument was later found to be essentially correct. In the language of Remark~\ref{NT:rem:classno}: for the
corresponding discriminants $D=-3,-4,-7,-8,-11,-19,-43,-67,-163$ there is exactly one reduced primitive form.
For $D=-3$ and $D=-4$ it is the hexagonal and the square lattice. For $D=-163$ it is $x^2+xy+41y^2$, and
Euler's observation that $k^2+k+41$ is prime for $0\le k\le39$ is a shadow of $h(-163)=1$ (Rabinowitsch, 1913).

\begin{theorem}[The $15$ theorem; Conway--Schneeberger 1993, Bhargava 2000]\label{NT:gal:15}
A positive definite quadratic form with integer Gram matrix represents every positive integer if and only if it
represents $1,2,3,5,6,7,10,14,15$.
\end{theorem}
Conway and Schneeberger did not publish their proof. M.~Bhargava, \emph{On the Conway--Schneeberger
fifteen theorem}, in \emph{Quadratic Forms and Their Applications} (Dublin 1999), Contemp.\ Math.\ 272,
AMS (2000), 27--37, gave a short proof, built by ``escalating'' lattices: one adds a vector whose
norm is the smallest integer not yet represented. This is the arithmetic version of the laminated
lattices of Lecture~6. None of the nine numbers can be omitted.

\begin{theorem}[The $290$ theorem; Bhargava--Hanke]\label{NT:gal:290}
A positive definite integer-valued quadratic form represents every positive integer if and only if it
represents
the $29$ numbers 1, 2, 3, 5, 6, 7, 10, 13, 14, 15, 17, 19, 21, 22, 23, 26, 29, 30, 31, 34, 35, 37, 42, 58, 93,
110, 145, 203, 290.
\end{theorem}
Conjectured by Conway. M.~Bhargava and J.~Hanke, \emph{Universal quadratic forms and the 290-theorem},
preprint (2011). It is widely cited, but as of 2026 it has not appeared in a journal.
``Integer-valued'' allows odd off-diagonal coefficients, such as $x^2+xy+y^2$; that is why the list is longer.

\begin{theorem}[Hardy 1916; Huxley 2003]\label{NT:gal:circle}
In the circle problem, $E(r)$ is not $o\big(r^{1/2}(\log r)^{1/4}\big)$. On the other hand,
$E(r)=O_\varepsilon(r^{131/208+\varepsilon})$ for every $\varepsilon>0$.
\end{theorem}
The lower bound is in G.~H.~Hardy, \emph{On Dirichlet's divisor problem}, Proc.\ London Math.\ Soc.\ (2)
15 (1916). The upper bound is in M.~N.~Huxley, \emph{Exponential sums and lattice points III}, Proc.\
London Math.\ Soc.\ (3) 87 (2003), 591--609. The logarithmic factors in the lower bound were sharpened
by Hafner (Invent.\ Math.\ 63, 1981) and by Soundararajan (Int.\ Math.\ Res.\ Not.\ 2003, no.~36). For
the arXiv improvements of the upper bound, see Section~\ref{NT:sec:circle}. This is the prototype of a boundary-layer error in a density count.

\begin{theorem}[Milnor 1964; Niemeier 1973]\label{NT:gal:niemeier}
There are non-isometric isospectral flat tori in dimension $16$. There are exactly $24$ even unimodular
lattices in $\R^{24}$, and the Leech lattice is the only one with no vectors of norm $2$.
\end{theorem}
J.~Milnor, \emph{Eigenvalues of the Laplace operator on certain manifolds}, Proc.\ Nat.\ Acad.\ Sci.\
USA 51 (1964), 542; H.-V.~Niemeier, J.\ Number Theory 5 (1973), 142--178; J.~H.~Conway, Invent.\ Math.\ 7
(1969), 137--142, for the characterization of Leech. Both are consequences of the theta series being a
modular form (Section~\ref{NT:sec:theta}).

\begin{theorem}[Cohn--Elkies 2003]\label{NT:gal:CE}
Theorem~\ref{NT:thm:CE} holds.
\end{theorem}
H.~Cohn and N.~Elkies, \emph{New upper bounds on sphere packings I}, Ann.\ of Math.\ 157 (2003), 689--714.
Their numerical bounds were the best known in many dimensions, and in dimensions $8$ and $24$ they came
extremely close to the densities of $E_8$ and Leech. This is the bridge from Blichfeldt's method to the solution in dimensions $8$ and $24$.

\begin{conjecture}[Cohn--Elkies]\label{NT:gal:CE2}
The linear programming bound of Theorem~\ref{NT:thm:CE} is sharp in dimension $2$, that is, its optimum
equals $\pi/\sqrt{12}$.
\end{conjecture}
Numerical optimization supports it, but no exact auxiliary function is known. Thue's
theorem is known by other means (Lecture~1), so what is open is whether \emph{this method} reproves it.

\begin{theorem}[Viazovska 2016; Cohn--Kumar--Miller--Radchenko--Viazovska 2016]\label{NT:gal:viazovska}
The $E_8$ lattice packing is the densest sphere packing in $\R^8$, and the Leech lattice packing is the
densest in $\R^{24}$.
\end{theorem}
M.~Viazovska, \emph{The sphere packing problem in dimension 8}, Ann.\ of Math.\ 185 (2017), 991--1015;
H.~Cohn, A.~Kumar, S.~D.~Miller, D.~Radchenko and M.~Viazovska, \emph{The sphere packing problem in dimension
24}, Ann.\ of Math.\ 185 (2017), 1017--1033. In both cases the magic function is built from (quasi)modular
forms, as sketched in Section~\ref{NT:sec:LP}.

\begin{theorem}[Universal optimality; Cohn--Kumar--Miller--Radchenko--Viazovska 2022]\label{NT:gal:universal}
Let $p(r)=g(r^2)$ with $g\colon(0,\infty)\to\R$ completely monotone. Among all point configurations in $\R^8$
(resp.\ $\R^{24}$) of density one point per unit volume, $E_8$ (resp.\ the Leech lattice) minimizes the
lower $p$-energy.
\end{theorem}
H.~Cohn, A.~Kumar, S.~D.~Miller, D.~Radchenko and M.~Viazovska, \emph{Universal optimality of the $E_8$ and
Leech lattices and interpolation formulas}, Ann.\ of Math.\ 196 (2022), no.~3, 983--1082. Sphere packing is
a limiting case (very steep potentials). This connects the lattice problem to the Boltzmann weights and
energy landscapes of Lectures~7--11: $E_8$ and Leech are ground states for a whole class of interactions,
and not only optimal hard spheres.

\begin{theorem}[Siegel 1945; Minkowski--Hlawka 1943]\label{NT:gal:siegel}
Theorem~\ref{NT:thm:siegel} and Corollary~\ref{NT:cor:MH} hold.
\end{theorem}
C.~L.~Siegel, \emph{A mean value theorem in geometry of numbers}, Ann.\ of Math.\ 46 (1945), 340--347;
E.~Hlawka, \emph{Zur Geometrie der Zahlen}, Math.\ Z.\ 49 (1943/44), 285--312. Minkowski had stated the bound
without a complete proof. The lattice analogue of the saturation bound of Lecture~1.

\begin{theorem}[Klartag 2025]\label{NT:gal:klartag}
There is a constant $c>0$ such that for all large $n$ there is a lattice packing in $\R^n$ of density at least
$c\,n^2\,2^{-n}$.
\end{theorem}
B.~Klartag, \emph{Lattice packing of spheres in high dimensions using a stochastically evolving ellipsoid},
Invent.\ Math.\ (2026), doi:10.1007/s00222-026-01412-w; arXiv:2504.05042 (2025). It is the first lower bound, for lattices or for general packings, that beats
$n\,2^{-n}$ by more than a logarithmic factor. The proof grows an ellipsoid that contains no nonzero point
of a lattice, by a random process. Compare the Minkowski--Hlawka bound
$2\zeta(n)2^{-n}$ above and the table in Lecture~3.

%======================================================================
\section*{Exercises}
\addcontentsline{toc}{section}{Exercises}

\begin{exercise}\label{NT:ex:reduce}
Reduce the form $26x^2+130xy+169y^2$ (the Gram matrix of the basis $(1,5)$, $(0,13)$ of the lattice in
Figure~\ref{NT:fig:fermat}). Check that you get $13(x^2+y^2)$, and read off a representation $13=a^2+b^2$.
Show that $h(-20)=2$ by listing the reduced forms of discriminant $-20$.
\end{exercise}

\begin{exercise}\label{NT:ex:hermite}
(Hermite's inequality.) Let $v_1$ be a shortest vector of $\Lat\subset\R^n$, and let $\Lat'$ be the
orthogonal projection of $\Lat$ onto $v_1^\perp$.
\begin{enumerate}[nosep]
\item Show that $\Lat'$ is a lattice with $\covol(\Lat')=\covol(\Lat)/\abs{v_1}$.
\item Show that every $w'\in\Lat'$ lifts to $w\in\Lat$ with $\abs{w}^2\le\abs{w'}^2+\frac14\abs{v_1}^2$, and
  deduce $\frac34\lambda_1(\Lat)^2\le\lambda_1(\Lat')^2$.
\item Conclude that $\gamma_n\le\big(\frac43\gamma_{n-1}\big)^{(n-1)/n}$, and hence $\gamma_n\le(4/3)^{(n-1)/2}$.
  Compare with the proof of Theorem~\ref{L1:thm:lattice2}.
\end{enumerate}
\end{exercise}

\begin{exercise}
Show that $N_\Lat(r)=V_nr^n/\covol(\Lat)+O(r^{n-1})$ for every lattice $\Lat\subset\R^n$, by Gauss's
argument with a fundamental parallelepiped. Then use it to show that the average of $r_2(m)$ for
$m\le X$ tends to $\pi$, and that the average of $r_4(m)$ over $m\le X$ is asymptotic to $\frac{\pi^2}{2}X$.
\end{exercise}

\begin{exercise}
Use Minkowski's theorem to show that every prime $p\equiv1$ or $3\pmod8$ can be written as $x^2+2y^2$.
(Use that $-2$ is a square mod $p$ for these primes, and the ellipse $x^2+2y^2<2p$, of area $\sqrt2\,\pi p>4p$.)
\end{exercise}

\begin{exercise}
Prove Corollary~\ref{NT:cor:thetatrans} for $\Lat=\Z$ directly, by applying Poisson summation to
$e^{-\pi tx^2}$, $t>0$. Deduce that $t^{1/4}\,\Theta_\Z(it)$ is invariant under $t\mapsto1/t$, and use this to
show that $\Theta_\Z(it)\sim t^{-1/2}$ as $t\to0^+$. Interpret this limit as a lattice point count.
\end{exercise}

\begin{exercise}\label{NT:ex:D16}
Show that $D_n^+=D_n\cup(D_n+\frac12\mathbf 1)$ is a lattice if and only if $n$ is even, that it is
unimodular when it is integral, and that it is even unimodular if and only if $8\mid n$. Count the vectors of
norm $2$ and $4$ in $D_{16}^+$ and compare with $E_8\oplus E_8$.
\end{exercise}

\begin{exercise}
Using $\Theta_{E_8}=E_4$, show that the number of vectors of norm $2m$ in $E_8$ is divisible by $240$ for every $m\ge1$.
Then verify by hand that $E_8$ has $2160$ vectors of norm $4$. (They are $\pm2e_i$, the vectors
$\pm e_i\pm e_j\pm e_k\pm e_l$, and the half-integral vectors with one coordinate $\pm\frac32$.)
\end{exercise}

\begin{exercise}
(The LP bound in dimension one.) Let $f(x)=\max(0,1-\abs x)$. Compute $\hat f(y)=\big(\frac{\sin\pi y}{\pi y}\big)^2$,
check that $f$ satisfies the hypotheses of Theorem~\ref{NT:thm:CE} with $r=1$, and conclude that the bound is
sharp in dimension $1$. Which lattice packings achieve it, and how does this show up in the equality analysis of
Section~\ref{NT:sec:LP}?
\end{exercise}

\begin{exercise}
Using Siegel's theorem, compute the average over $X_n$ of the number of nonzero lattice vectors in a ball
of volume $V$ centred at $0$. Deduce that some unimodular lattice has $\lambda_1\ge(2/V_n)^{1/n}$, which is
asymptotically $\sqrt{n/(2\pi e)}$. Compare with the upper bound $\lambda_1\le2V_n^{-1/n}$ coming from
Minkowski's theorem, and translate both into bounds on $\gamma_n$.
\end{exercise}

\section*{Further reading}
\addcontentsline{toc}{section}{Further reading}
\begin{itemize}[nosep]
\item J.-P.~Serre, \emph{A Course in Arithmetic}, Springer GTM 7, 1973. Chapter~V: quadratic forms over $\Z$
  (including $8\mid n$ for even unimodular lattices); Chapter~VII: modular forms and theta series.
\item J.~H.~Conway and N.~J.~A.~Sloane, \emph{Sphere Packings, Lattices and Groups}, 3rd ed., Springer, 1999.
  Ch.~2 (theta series), Ch.~16--18 (Niemeier lattices, Leech).
\item J.~W.~S.~Cassels, \emph{An Introduction to the Geometry of Numbers}, Springer, 1959. Blichfeldt,
  Minkowski, Siegel and Hlawka.
\item J.~Martinet, \emph{Perfect Lattices in Euclidean Spaces}, Springer, 2003. Hermite constants and Voronoi's theory.
\item J.~H.~Conway (with F.~Fung), \emph{The Sensual (Quadratic) Form}, MAA, 1997. Binary forms
  and the topograph, and the 15 theorem.
\item H.~Cohn, \emph{A conceptual breakthrough in sphere packing}, Notices AMS 64 (2017), 102--115. The LP
  bound and Viazovska's construction, for non-specialists.
\item M.~N.~Huxley, \emph{Area, Lattice Points, and Exponential Sums}, Oxford, 1996. For the circle problem.
\end{itemize}
